\documentclass[11pt,a4paper]{article}
\usepackage[margin=22mm]{geometry}
\usepackage[T1]{fontenc}
\usepackage[utf8]{inputenc}
\usepackage[english]{babel}
\usepackage{amsmath,amssymb,mathtools}
\usepackage{booktabs,longtable,array}
\usepackage{xcolor,listings}
\usepackage[bookmarks=false]{hyperref}
\hypersetup{colorlinks=true,linkcolor=blue,urlcolor=blue}
\newcommand{\ket}[1]{\lvert #1\rangle}
\newcommand{\bra}[1]{\langle #1\rvert}
\newcommand{\code}[1]{\nolinkurl{#1}}
\newcommand{\Tr}{\operatorname{Tr}}
\newcommand{\ii}{\mathrm i}
\newcommand{\Bf}{\mathcal B_f}
\lstset{basicstyle=\ttfamily\scriptsize,numbers=left,numberstyle=\tiny\color{gray},
  numbersep=7pt,frame=single,breaklines=true,columns=fullflexible,showstringspaces=false,
  keepspaces=true,upquote=true}
\emergencystretch=3em
\title{Source-to-equation guide to the projected-Fock scattering code\\
\large Every computational line, in execution order}
\author{Documentation for \code{codex/code}}
\date{Source inspected 25 September 2026}
\begin{document}
\maketitle
\begin{abstract}
This is a reading guide to the actual source. Each listing is copied from the current
Python files; the table after it accounts for every executable statement. One table
row can cover the continuation lines of a single Python statement. Comments and
blank lines have no operation. Equations specify the array or matrix entries produced.
The notebook section describes its current cells, including settings that differ
from the function defaults.
\end{abstract}
\tableofcontents

\section{The index convention used everywhere}
\begin{center}\small
\begin{tabular}{p{.20\textwidth}p{.69\textwidth}}
\toprule Symbol / Python object & Exact meaning\\\midrule
$M=\texttt{basis.modes}$ & Number of retained signed field momenta. Python mode index $m=0,\ldots,M-1$.\\
$N=\texttt{basis.cap}$ & Photon occupation cap; its constraint depends on \code{truncation}.\\
$k_m=\texttt{k\_tab[m]}$ & Signed wavevector; $\omega_m=|k_m|$ because the code uses $c=1$.\\
$\mathbf n=(n_0,\ldots,n_{M-1})$ & Field occupation tuple; every tuple has exactly $M$ entries.\\
$i$ & \code{basis.index[occupation]}: position of a tuple in \code{basis.states}; not a mode index.\\
$s\in\{0,1\}$ & TLS index: $s=0$ is $\ket g$, $s=1$ is $\ket e$.\\
$q=2i+s$ & Position of $\ket{\mathbf n_i;s}$ in a ket or Hamiltonian matrix.\\
$v_q=C_{i,s}$ & Complex amplitude in $\ket\psi=\sum_{i,s}C_{i,s}\ket{\mathbf n_i;s}$.\\
$B=|\Bf|,\ d=2B$ & Number of occupation tuples; ket length and Hamiltonian dimension.\\
$L,D,\omega_0,x_a$ & Box length, coupling parameter, TLS frequency, TLS position; keys \code{L}, \code{D}, \code{omega_0}, \code{x_tls}.\\
$t_j$ & Requested output time. Index $j$ enumerates saved states, not internal ODE steps.\\
\bottomrule
\end{tabular}
\end{center}
The indexing identity required by almost every file is
\begin{equation}
\boxed{v[2i]=C_{i,g},\qquad v[2i+1]=C_{i,e},\qquad
v[2i:2i+2]=(C_{i,g},C_{i,e}).}
\label{eq:slice}
\end{equation}
Python excludes the right endpoint of a slice, so \code{v[2*i:2*i+2]}
selects exactly two cells. Similarly, \code{v[0::2]} selects every even
cell and \code{v[1::2]} every odd cell.

\subsection{Execution graph and shapes}
\begin{align*}
\texttt{four dictionaries}
&\xrightarrow{\texttt{ExperimentConfig}}
\texttt{Experiment}\\
&\xrightarrow{\texttt{momentum\_modes}}k\in\mathbb R^M
\xrightarrow{\texttt{FockBasis}}\{\mathbf n_i\}_{i=0}^{B-1}\\
&\xrightarrow{\texttt{Hamiltonian}}H_0,V\in\mathbb C^{d\times d}
\xrightarrow{\texttt{initial\_state}}v_0\in\mathbb C^d\\
&\xrightarrow{\texttt{sesolve}}v(t_j)\in\mathbb C^d
\xrightarrow{\texttt{even/odd slices}}C_g,C_e\in\mathbb C^{J\times B}\\
&\xrightarrow{\texttt{compute\_observables}}
\mathcal N,P_e,T_1,R_1,P_n\in\mathbb R^J.
\end{align*}
Here $J=\lfloor T/\Delta t\rfloor+1$ if all states are stored. The program
propagates a QuTiP ket; its $d\times d$ Hamiltonian is sparse. A global
density matrix is never propagated.

\section{Configuration: \code{src/xp_config.py}}
\lstinputlisting[firstline=1,lastline=15,firstnumber=1]{../src/xp_config.py}
\begin{longtable}{p{.14\textwidth}p{.76\textwidth}}
\toprule Lines & Exact role\\\midrule\endfirsthead
\toprule Lines & Exact role\\\midrule\endhead
1 & Import \code{dataclass}, which generates an initializer assigning declared fields.\\
4--5 & Make \code{ExperimentConfig} such a data container.\\
6 & Docstring only; no computation.\\
8 & \code{param_photon}: $k_0,\sigma_k,x_0$ and optional field-preparation switches.\\
9 & \code{param_atom}: $\omega_0,D,x_a,L$, coupling profile, optional TLS preparation.\\
10 & \code{param_time_evol}: $T,\Delta t$ and optional solver controls.\\
11 & \code{cutoffs}: IR/UV bounds or explicit signed \code{k_modes}.\\
12 & Default photon cap $N=3$.\\
13 & Default basis $\sum_mn_m\leq N$.\\
14 & Default \code{False} retains counter-rotating transitions.\\
15 & Default \code{True} requests the complete saved ket trajectory.\\
\bottomrule
\end{longtable}
No grid, state or Hamiltonian is built in this class.

\section{Basis construction: \code{src/states.py:10--36}}
\lstinputlisting[firstline=10,lastline=36,firstnumber=10]{../src/states.py}
\begin{longtable}{p{.14\textwidth}p{.76\textwidth}}
\toprule Lines & Exact operation / equation\\\midrule\endfirsthead
\toprule Lines & Exact operation / equation\\\midrule\endhead
10--11 & \code{FockBasis(modes,cap,truncation)} receives $M,N$ and the basis rule.\\
12--14 & Store the three inputs on the object.\\
15 & \code{(0,)*modes} repeats zero $M$ times: $\mathbf0=(0,\ldots,0)$.\\
16--17 & Start the one-occupied-mode basis with exactly one field vacuum.\\
18--19 & Iterate $m=0,\ldots,M-1$ and $n=1,\ldots,N$: exactly $MN$ pairs.\\
20--22 & Copy vacuum, set entry $m$ to $n$, append tuple $n\mathbf e_m$. Thus $\Bf^{\rm one}=\{\mathbf0\}\cup\{n\mathbf e_m\}$ and $|\Bf^{\rm one}|=1+MN$. A tuple is hashable for the dictionary at line 34.\\
23--25 & For total cap, iterate total photon count $r=0,\ldots,N$.\\
26 & \code{combinations_with_replacement(range(M),r)} lists multisets of $r$ mode indices; a repeated index denotes multiple photons in that mode.\\
27 & \code{np.bincount(...,minlength=M)} converts the multiset into $(n_0,\ldots,n_{M-1})$, $\sum_mn_m=r$; append this tuple once.\\
28--29 & \code{product(range(N+1),repeat=M)} enumerates $\{0,\ldots,N\}^M$: an independent cap in each mode.\\
30--31 & Reject any unknown basis name.\\
33 & Store the selected field tuples as \code{self.states}.\\
34 & Create inverse map $\texttt{index}[\mathbf n_i]=i$. \code{enumerate} supplies $i$; later matrix construction uses this same map.\\
35 & $\texttt{photon\_numbers}[i]=\sum_{m=0}^{M-1}n_m(i)$.\\
36 & Two TLS states per tuple: $d=2|\Bf|$.\\
\bottomrule
\end{longtable}
\subsection{The three dimensions}
\begin{align}
d_{\rm one}&=2(1+MN),\\
d_{\rm per}&=2(N+1)^M,\\
d_{\rm tot}&=2\sum_{r=0}^{N}\binom{M+r-1}{r}
             =2\binom{M+N}{N}.
\end{align}
For the last count, distributing $r$ indistinguishable photons among $M$ modes
gives $\binom{M+r-1}{r}$. Alternatively add
$n_{\rm slack}=N-\sum_mn_m\geq0$ and count nonnegative solutions of
$n_0+\cdots+n_{M-1}+n_{\rm slack}=N$.
For $M=2,N=2$:
\begin{center}\small
\begin{tabular}{p{.22\textwidth}p{.54\textwidth}r}
\toprule Rule & Field occupations & $d$\\\midrule
\code{truncated} & $(0,0),(1,0),(2,0),(0,1),(0,2)$ & 10\\
\code{full+totalcap} & preceding five, plus $(1,1)$ & 12\\
\code{full} & all $(n_0,n_1)\in\{0,1,2\}^2$ & 18\\
\bottomrule
\end{tabular}
\end{center}
\section{Initial ket: \code{src/states.py:39--93}}
\subsection{Packet and TLS coefficients}
\lstinputlisting[firstline=39,lastline=64,firstnumber=39]{../src/states.py}
\begin{equation}
\widetilde c_m=
e^{-(k_m-k_0)^2/(4\sigma_k^2)}
e^{-\ii k_m x_0}\mathbf1_{k_m>0}
\quad\text{by default},\qquad
c_m=\frac{\widetilde c_m}{\sqrt{\sum_j|\widetilde c_j|^2}}.
\label{eq:packet}
\end{equation}
\begin{longtable}{p{.14\textwidth}p{.76\textwidth}}
\toprule Lines & Exact operation / equation\\\midrule\endfirsthead
\toprule Lines & Exact operation / equation\\\midrule\endhead
39--40 & Function receives basis, signed $k$ array, photon and TLS dictionaries; docstring has no operation.\\
41--42 & Read $k_0$ and $\sigma_k$.\\
43--44 & Require $k_0>0,\sigma_k>0$ for this packet convention.\\
46 & Real Gaussian amplitude $e^{-(k_m-k_0)^2/(4\sigma_k^2)}$ for every $m$. Its squared modulus has variance $\sigma_k^2$ before discrete normalization.\\
47 & Multiply by $e^{-\ii k_mx_0}$; the phase moves $\sum_m c_me^{\ii k_mx}$ toward $x_0$. The array becomes complex.\\
48--49 & Default \code{right_moving_only=True}: zero all $k_m\leq0$. With \code{False}, remove the indicator in Eq.~\eqref{eq:packet}.\\
50--52 & Compute $\|\widetilde c\|_2$ and reject zero support, preventing division by zero.\\
53 & In-place normalization gives $\sum_m|c_m|^2=1$; subsequent \code{packet[m]} means $c_m$.\\
55 & Read TLS preparation; absent key means \code{g}.\\
56--61 & Map \code{g}, \code{e}, \code{+}, \code{-} to $\eta=(\eta_g,\eta_e)$: $(1,0)$, $(0,1)$, $(1,1)/\sqrt2$, $(1,-1)/\sqrt2$.\\
62--63 & Reject other TLS labels.\\
64 & Convert $\eta$ to a two-entry complex NumPy array.\\
\bottomrule
\end{longtable}

\subsection{Number state: why \code{v[2*i:2*i+2] += ...}}
\lstinputlisting[firstline=66,lastline=79,firstnumber=66]{../src/states.py}
\begin{longtable}{p{.14\textwidth}p{.76\textwidth}}
\toprule Lines & Exact operation / equation\\\midrule\endfirsthead
\toprule Lines & Exact operation / equation\\\midrule\endhead
66 & Absent \code{state} means \code{number}.\\
67 & Allocate $v\in\mathbb C^d$, $v_q=0$ for all $q$. Complex dtype retains phases.\\
68--72 & Select number branch; read integer $n$ (default 1), require $1\leq n\leq N$.\\
74 & \code{enumerate(packet)} produces $(m,c_m)$ for every retained mode.\\
75--76 & Construct $\mathbf n^{(m)}=n\mathbf e_m$: all modes empty except $m$.\\
77 & Look up $i_m=\texttt{basis.index.get}(\mathbf n^{(m)})$; missing tuple returns \code{None}.\\
78 & Skip absent occupation. \code{is not None} is required because $i_m=0$ is a valid index. Current three bases include every $n\mathbf e_m$ for $1\leq n\leq N$.\\
79 & By Eq.~\eqref{eq:slice}, add $c_m(\eta_g,\eta_e)$ to the two amplitudes belonging to $\mathbf n^{(m)}$:
$C_{i_m,g}\mathrel{+}=c_m\eta_g$, $C_{i_m,e}\mathrel{+}=c_m\eta_e$.
After the loop,
$\displaystyle\ket{\psi(0)}
=\left[\sum_m c_m\ket{n\mathbf e_m}\right]\otimes
(\eta_g\ket g+\eta_e\ket e)$.
\code{+=} accumulates contributions; for $n\geq1$ the occupations in this loop are distinct.\\
\bottomrule
\end{longtable}
For default $n=1,\eta=(1,0)$, line 79 writes
$v_{2i_m}=c_m$, $v_{2i_m+1}=0$ and
$\ket{\psi(0)}=\sum_m c_m\ket{1_m;g}$.
For $n>1$ the code uses $\sum_m c_m\ket{n_m}$, not
$(\sum_m c_ma_m^\dagger)^n\ket0/\sqrt{n!}$; the latter includes cross-mode terms.

\subsection{Coherent state and projection}
\lstinputlisting[firstline=80,lastline=93,firstnumber=80]{../src/states.py}
Set $\alpha_m=\alpha c_m$. The unprojected product state has coefficient
\begin{equation}
\langle\mathbf n|\boldsymbol\alpha\rangle
=e^{-|\alpha|^2/2}
\prod_{m=0}^{M-1}\frac{(\alpha c_m)^{n_m}}{\sqrt{n_m!}},
\qquad \sum_m|c_m|^2=1.
\label{eq:coherent}
\end{equation}
\begin{longtable}{p{.14\textwidth}p{.76\textwidth}}
\toprule Lines & Exact operation / equation\\\midrule\endfirsthead
\toprule Lines & Exact operation / equation\\\midrule\endhead
80--81 & Select coherent branch; read \code{alpha} (default 1), cast to complex.\\
82 & Visit every retained occupation $\mathbf n_i$ and its index $i$. This loop implements projection by never visiting omitted tuples.\\
83 & Initialize field coefficient at $e^{-|\alpha|^2/2}=\prod_m e^{-|\alpha c_m|^2/2}$.\\
84--85 & Multiply by $(\alpha c_m)^{n_m}/\sqrt{n_m!}$ for each field mode. Result is Eq.~\eqref{eq:coherent}.\\
86 & Assign $(C_{i,g},C_{i,e})=\texttt{coefficient}\,(\eta_g,\eta_e)$; each tuple occurs once.\\
87--88 & Reject unknown preparation name.\\
90--92 & Compute $\|v\|_2$ and reject zero support. Projection can make this norm less than one.\\
93 & Return QuTiP ket $v/\|v\|_2=P_{\Bf}\ket{\boldsymbol\alpha,\eta}/\|P_{\Bf}\ket{\boldsymbol\alpha,\eta}\|$.\\
\bottomrule
\end{longtable}
Different bases can therefore produce different \emph{normalized} coherent
initial states; their fidelity may already be below one at $t=0$.

\section{Hamiltonian assembly: \code{src/hamiltonians.py}}
\subsection{Mathematical target and constructor}
Write $f_m=1$ for \code{flat} and $f_m=\sqrt{|k_m|}$ for \code{sqrt}.
The finite Hamiltonian constructed by the code is
\begin{align}
H(D)&=H_0+D V_1,\label{eq:HD}\\
H_0&=\sum_m|k_m|a_m^\dagger a_m
       +\omega_0\ket e\bra e,\label{eq:H0}\\
V_1&=\sum_m\left(u_ma_m+u_m^*a_m^\dagger\right)
                  (\sigma^++\sigma^-),\qquad
u_m=\frac{\ii f_m e^{-\ii k_mx_a}}{\sqrt L}.
\label{eq:V}
\end{align}
Here $\sigma^+=\ket e\bra g$, $\sigma^-=\ket g\bra e$.
The actual operator is projected onto the enumerated finite basis:
$P_{\mathcal B}H(D)P_{\mathcal B}$.
\lstinputlisting[firstline=8,lastline=15,firstnumber=8]{../src/hamiltonians.py}
\begin{longtable}{p{.14\textwidth}p{.76\textwidth}}
\toprule Lines & Exact operation\\\midrule\endfirsthead
\toprule Lines & Exact operation\\\midrule\endhead
8--9 & Constructor receives basis, signed momentum array, TLS dictionary and RWA switch.\\
10 & Store the very same \code{FockBasis} whose indices define ket amplitudes.\\
11 & Convert \code{k_tab} to a NumPy array so multiplication and absolute value act elementwise.\\
12--13 & Store physical parameters and RWA flag.\\
14 & Call \code{free()} immediately; store sparse diagonal $H_0$.\\
15 & Call \code{interaction()} immediately; store sparse unit-coupling $V_1$. This construction runs before \code{sesolve}.\\
\bottomrule
\end{longtable}

\subsection{Free diagonal: every assigned entry}
\lstinputlisting[firstline=17,lastline=24,firstnumber=17]{../src/hamiltonians.py}
For occupation tuple $\mathbf n_i$ define
$E_f(i)=\sum_{m=0}^{M-1}|k_m|n_m(i)$.
Then
\begin{equation}
(H_0)_{2i,2i}=E_f(i),\qquad
(H_0)_{2i+1,2i+1}=E_f(i)+\omega_0.
\label{eq:Hdiag}
\end{equation}
\begin{longtable}{p{.14\textwidth}p{.76\textwidth}}
\toprule Lines & Exact operation\\\midrule\endfirsthead
\toprule Lines & Exact operation\\\midrule\endhead
17--18 & Enter method and read $\omega_0$.\\
19 & Allocate real array \code{energies} of length $d$, initially zero. No dense $d\times d$ matrix is allocated.\\
20 & Iterate each occupation tuple $\mathbf n_i$ and its index $i$.\\
21 & \code{np.abs(k_tab)} is $(|k_0|,\ldots,|k_{M-1}|)$; dot product with the tuple is $E_f(i)$.\\
22 & Position $2i$ is $\ket{\mathbf n_i;g}$, giving first diagonal in Eq.~\eqref{eq:Hdiag}.\\
23 & Position $2i+1$ is $\ket{\mathbf n_i;e}$; add $\omega_0$ for TLS excitation.\\
24 & \code{diags(...,format="csr")} creates compressed-sparse-row diagonal $H_0$.\\
\bottomrule
\end{longtable}

\subsection{Profile and every sparse transition}
\lstinputlisting[firstline=26,lastline=60,firstnumber=26]{../src/hamiltonians.py}
Bosonic matrix elements are
\begin{equation}
a_m\ket{\mathbf n;s}=\sqrt{n_m}\ket{\mathbf n-\mathbf e_m;s},
\quad
a_m^\dagger\ket{\mathbf n;s}=\sqrt{n_m+1}\ket{\mathbf n+\mathbf e_m;s}.
\label{eq:boson}
\end{equation}
The TLS factor flips $s\mapsto1-s$. Consequently, with source column
$q=2i+s$ and target row $p=2j+(1-s)$,
\begin{align}
V_{1,pq}^{(-)}&=u_m\sqrt{n_m}
&&\text{if }\mathbf n_j=\mathbf n_i-\mathbf e_m,\label{eq:down}\\
V_{1,pq}^{(+)}&=u_m^*\sqrt{n_m+1}
&&\text{if }\mathbf n_j=\mathbf n_i+\mathbf e_m.\label{eq:up}
\end{align}
The matrix entry is $H_{pq}=D V_{1,pq}$; matrix convention is
$H_{pq}=\bra{\text{target }p}H\ket{\text{source }q}$.
\begin{longtable}{p{.14\textwidth}p{.76\textwidth}}
\toprule Lines & Exact operation / equation\\\midrule\endfirsthead
\toprule Lines & Exact operation / equation\\\midrule\endhead
26--29 & Enter interaction builder; read $L,x_a$ and profile name.\\
30--31 & For \code{sqrt}, vector $f_m=\sqrt{|k_m|}$; each nonzero mode gets frequency-dependent coupling.\\
32--33 & For \code{flat}, vector $f_m=1$.\\
34--35 & Reject any other coupling profile.\\
36 & Vectorized formula $u_m=\ii f_m e^{-\ii k_mx_a}/\sqrt L$. This is \emph{unit} $D$; line 65 inserts $D$. \code{1j} is Python's $\ii$.\\
38 & Three empty lists will hold target rows $p$, source columns $q$, and complex entries $V_{1,pq}$.\\
39 & Loop over every valid source occupation $\mathbf n_i$. This visits $B$ tuples, not $B^2$ pairs.\\
40 & Loop over source TLS index $s=0,1$.\\
41 & Compute source ket index $q=2i+s$ from Eq.~\eqref{eq:slice}.\\
42 & Try a transition in each field mode $m=0,\ldots,M-1$.\\
43 & Downward transition requires $n_m>0$. Without RWA allow either $s$. With RWA allow only $s=0$: annihilate a photon and excite TLS, $a_m\sigma^+$.\\
44--45 & Copy the occupation tuple and subtract one in mode $m$, giving $\mathbf n-\mathbf e_m$ while preserving all other fields.\\
46 & Look up target occupation index $j$. Direct dictionary access is safe because all three bases are downward closed: removing one photon from any retained tuple remains retained.\\
47 & Append target row $p=2j+1-s$; $1-s$ flips $g\leftrightarrow e$.\\
48 & Append source column $q$.\\
49 & Append $u_m\sqrt{n_m}$, exactly Eq.~\eqref{eq:down}; the square root is required by $a_m\ket{n_m}=\sqrt{n_m}\ket{n_m-1}$.\\
50 & Upward transition: without RWA allow either $s$; with RWA allow only $s=1$ so TLS de-excites while a photon is created, $a_m^\dagger\sigma^-$.\\
51--52 & Copy tuple and increment field mode $m$, giving $\mathbf n+\mathbf e_m$.\\
53 & \code{index.get} returns \code{None} when the new occupation violates the chosen truncation. This realizes $P_{\mathcal B}VP_{\mathcal B}$ at the upper boundary.\\
54 & Include upward edge only if target is in the basis. Again $j=0$ is valid; testing \code{j is not None} handles it.\\
55--56 & Append target row $2j+1-s$ and source column $q$.\\
57 & Append $u_m^*\sqrt{n_m+1}$, Eq.~\eqref{eq:up}; complex conjugation is required for Hermiticity of $u_ma_m+u_m^*a_m^\dagger$.\\
59--60 & Construct COO sparse matrix from triplets $(p,q,V_{1,pq})$ with shape $d\times d$, then convert to CSR for matrix-vector operations. A COO entry at \code{rows[r],cols[r]} gets \code{values[r]}.\\
\bottomrule
\end{longtable}
For one mode, $N=1$, $x_a=0$, and $f=\sqrt{k}$, the order
$(\ket{0;g},\ket{0;e},\ket{1;g},\ket{1;e})$ gives
\[
H=\begin{pmatrix}
0&0&0&D u\\
0&\omega_0&D u&0\\
0&D u^*&k&0\\
D u^*&0&0&k+\omega_0
\end{pmatrix},
\qquad u=\frac{\ii\sqrt{k}}{\sqrt L}.
\]
Rows and columns are displayed to expose the index rule. With RWA, only
the $(\ket{0;e},\ket{1;g})$ block remains coupled.
For a nonzero $x_a$, $u$ carries the phase in Eq.~\eqref{eq:V}.

\subsection{Adding the physical coupling}
\lstinputlisting[firstline=62,lastline=65,firstnumber=62]{../src/hamiltonians.py}
\begin{longtable}{p{.14\textwidth}p{.76\textwidth}}
\toprule Lines & Exact operation\\\midrule\endfirsthead
\toprule Lines & Exact operation\\\midrule\endhead
62 & \code{D=None} allows either configured coupling or an explicit override.\\
63--64 & If no override, read $D=\texttt{param\_atom["D"]}$.\\
65 & Form $H(D)=H_0+D V_1$ as a sparse sum and wrap it as a QuTiP \code{Qobj}. No global density matrix is created.\\
\bottomrule
\end{longtable}

\section{Grid and resource count: \code{src/experiment.py:12--57}}
\subsection{Signed momenta}
\lstinputlisting[firstline=12,lastline=25,firstnumber=12]{../src/experiment.py}
For periodic length $L$, the candidate momenta are $k_j=2\pi j/L$.
The automatic branch returns
\begin{equation}
\mathcal K=\left\{\frac{2\pi j}{L}:
|j|\leq j_{\max},\quad
\Omega_{\rm IR}-10^{-12}\leq\left|\frac{2\pi j}{L}\right|
\leq\Omega_{\rm UV}+10^{-12}\right\}.
\label{eq:grid}
\end{equation}
\begin{longtable}{p{.14\textwidth}p{.76\textwidth}}
\toprule Lines & Exact operation\\\midrule\endfirsthead
\toprule Lines & Exact operation\\\midrule\endhead
12--13 & Function takes atom parameters and cutoff dictionary; docstring states signed periodic momenta.\\
14 & If \code{cutoffs} contains \code{k_modes}, use those prescribed values instead of generating an IR/UV grid.\\
15 & Cast explicit values to one NumPy float array. Their phases and $|k|$ later use these exact numbers.\\
16--17 & Require one-dimensional, nonempty and pairwise distinct wavevectors. \code{np.unique} counts exact distinct floats.\\
18 & Sort explicit $k$ from negative to positive; source-file mode index $m$ follows this sorted order.\\
20 & Read $L$.\\
21 & $j_{\max}=\lceil\Omega_{\rm UV}L/(2\pi)\rceil$. Ceiling ensures the candidate integer range reaches the UV edge.\\
22 & \code{np.arange(-m_max,m_max+1)} gives all integers from $-j_{\max}$ through $+j_{\max}$, inclusive; multiply by $2\pi/L$.\\
23--24 & Boolean mask is the intersection of lower and upper inequalities in Eq.~\eqref{eq:grid}. \code{&} is elementwise logical AND. The $10^{-12}$ margin protects an intended boundary from floating-point roundoff.\\
25 & Boolean indexing \code{k_tab[band]} returns only retained momenta. If $\Omega_{\rm IR}>0$, zero is removed automatically.\\
\bottomrule
\end{longtable}
In the example $L=20\pi$, $\Omega_{\rm IR}=0.8$,
$\Omega_{\rm UV}=1.2$, the spacing is $\Delta k=0.1$ and $M=10$:
\[
\mathcal K=\{-1.2,-1.1,-1.0,-0.9,-0.8\}
\cup\{0.8,0.9,1.0,1.1,1.2\}.
\]
An explicit \code{k_modes} array is not checked against periodic quantization
or for zero; that branch implements the supplied finite set literally.

\subsection{Dimension formulas and memory estimate}
\lstinputlisting[firstline=28,lastline=57,firstnumber=28]{../src/experiment.py}
\begin{longtable}{p{.14\textwidth}p{.76\textwidth}}
\toprule Lines & Exact operation / equation\\\midrule\endfirsthead
\toprule Lines & Exact operation / equation\\\midrule\endhead
28--31 & Signature and docstring. \code{param_photon} is absent because basis size and saved-state memory do not depend on packet coefficients.\\
33 & $M=|\mathcal K|$, computed by the same grid function used in the experiment.\\
34 & $N=\texttt{n\_max}$.\\
35--36 & Reject $M=0$ or $N<1$; a photon packet needs at least one mode.\\
37--41 & Store three exact ket dimensions derived in the basis section: $2\binom{M+N}{N}$, $2(1+MN)$, and $2(N+1)^M$. \code{comb} computes an integer binomial coefficient.\\
42--43 & Reject a name not in the dimension dictionary.\\
45 & Read requested final time $T$ and output spacing $\Delta t$.\\
46--47 & Require positive $T,\Delta t$.\\
49 & Number of output samples $J=\lfloor T/\Delta t\rfloor+1$, corresponding to $t_j=j\Delta t$, $j=0,\ldots,J-1$. If $T/\Delta t$ is not integer, $t_{J-1}<T$.\\
50 & Choose $d$ using the basis name.\\
51 & Complex128 amplitude is 16 bytes. Approximate ket-history storage in GiB is $16dJ/2^{30}$ if all states stored, or $16d/2^{30}$ if only final ket. This excludes sparse $H$, list/QuTiP overhead, NumPy copies, and observable arrays.\\
52 & Advisory threshold: \code{feasible} iff $d\le100\,000$ and estimated history $\le1$ GiB. This function only returns the Boolean; \code{Experiment} does \emph{not} enforce it.\\
53--55 & Print $M$, scheme, $d$, $J$, estimated GiB and advisory Boolean. The \code{f"..."} expressions format values; \code{:,} inserts thousands separators, \code{:.3g} prints three significant digits.\\
56--57 & Return these five numbers/Boolean in a Python dictionary for a notebook to inspect.\\
\bottomrule
\end{longtable}
For the current notebook parameters $M=10,N=3$,
\[
d_{\rm full}=2(4)^{10}=2{,}097{,}152,\quad
J=281,\quad
16dJ/2^{30}\simeq8.78\ {\rm GiB}.
\]
The notebook presently prints \code{feasible=False} but still starts that run;
the call site is analyzed below.

\section{\code{Experiment}: construction and propagation}
\subsection{Object construction, \code{src/experiment.py:60--85}}
\lstinputlisting[firstline=60,lastline=85,firstnumber=60]{../src/experiment.py}
\begin{longtable}{p{.14\textwidth}p{.76\textwidth}}
\toprule Lines & Exact operation / dependency\\\midrule\endfirsthead
\toprule Lines & Exact operation / dependency\\\midrule\endhead
60--63 & Class and constructor; input is one \code{ExperimentConfig}.\\
64--66 & Keep direct references to the three input dictionaries; the object can read their values later.\\
67--70 & Copy $N$, scheme, RWA switch, and storage switch from config.\\
72 & Build $\mathcal K=(k_0,\ldots,k_{M-1})$ using Eq.~\eqref{eq:grid}.\\
73 & $M=\texttt{len(k\_tab)}$; subsequent arrays use this \emph{actual} mode count.\\
75--76 & Reject a band with zero modes before trying to construct packet or Hamiltonian.\\
78 & Construct occupation list, inverse index, photon-count array and ket dimension as described above.\\
79 & Alias the same photon-count array for later sector masks.\\
80 & Construct $H_0$ and $V_1$ immediately. This is where sparse matrix assembly cost occurs.\\
81 & Form QuTiP $H=H_0+D V_1$.\\
82 & Construct and normalize initial QuTiP ket with the \emph{same} basis ordering and $k$ array.\\
83--85 & Set future coefficient arrays and observables to \code{None}; they are unavailable until propagation and measurement methods run.\\
\bottomrule
\end{longtable}
No solver is run in \code{__init__}; however, allocating \code{FockBasis}
and building $H$ can already be expensive. The memory Boolean above does not
stop this constructor.

\subsection{QuTiP call and every output slice, \code{src/experiment.py:87--110}}
\lstinputlisting[firstline=87,lastline=110,firstnumber=87]{../src/experiment.py}
\begin{equation}
\ii\,\partial_t v(t)=H v(t),\qquad v(0)=v_0,\qquad
t_j=j\Delta t,\quad
j=0,\ldots,\left\lfloor T/\Delta t\right\rfloor.
\label{eq:ode}
\end{equation}
\begin{longtable}{p{.14\textwidth}p{.76\textwidth}}
\toprule Lines & Exact operation / equation\\\midrule\endfirsthead
\toprule Lines & Exact operation / equation\\\midrule\endhead
87--88 & Method and docstring; \code{progress=False} disables progress display unless requested.\\
89 & Read $T,\Delta t$ from dictionary.\\
90--91 & Require positive values.\\
92 & \code{np.arange(J)*dt} creates $t_j=j\Delta t$. This is an \emph{output} grid; BDF chooses internal steps independently. Last saved time is $\lfloor T/\Delta t\rfloor\Delta t$.\\
93 & Solver method from dictionary, default \code{"bdf"}.\\
94--95 & \code{qutip.sesolve(H,state0,times,...)} numerically solves Eq.~\eqref{eq:ode}; input is a ket, so state dimension is $d$ rather than $d^2$.\\
96 & Pass selected integrator method.\\
97 & If \code{store_state=True}, store every ket $v(t_j)$; otherwise retain no trajectory list.\\
98 & Request final ket in either storage mode.\\
99 & \code{normalize_output=False}: do not force $\|v(t_j)\|=1$ after numerical integration. Norm deviation remains observable.\\
100--101 & Relative/absolute ODE tolerances, defaults $10^{-9},10^{-11}$. These constrain local integration error according to solver conventions, not the Fock/mode truncation error.\\
102 & Show \code{tqdm} progress bar only if \code{progress=True}.\\
103 & End the solver call; \code{self.result} contains QuTiP states and metadata.\\
104--105 & If full storage, convert every QuTiP ket to a one-dimensional complex NumPy vector via \code{state.full()[:,0]}; stack rows to shape $(J,d)$. \code{[:,0]} removes the single-column dimension of a ket.\\
106 & \code{vectors[:,0::2]} extracts $C_g(t_j,i)$ and \code{vectors[:,1::2]} extracts $C_e(t_j,i)$, both shape $(J,B)$ by Eq.~\eqref{eq:slice}.\\
107--108 & Final-only branch extracts one vector of shape $(d,)$.\\
109 & Even/odd slices now yield two arrays of shape $(B,)$; no time axis.\\
110 & Return the arrays. They are also stored as object attributes.\\
\bottomrule
\end{longtable}
QuTiP's solver errors and finite-model errors are separate:
\[
\underbrace{\|v_{\rm numerical}(t)-v_{\rm exact\ finite}(t)\|}_{\text{ODE error}}
\quad\text{versus}\quad
\underbrace{\|v_{\rm exact\ finite}(t)-\psi_{\rm target}(t)\|}_{\text{basis/grid model error}}.
\]
The saved norm diagnoses the first on the chosen finite Hamiltonian; it cannot
certify that omitted Fock states or momentum modes are unimportant.

\section{Scattering observables: \code{src/experiment.py:112--152}}
\subsection{Selecting the one-waveguide-photon sector}
\lstinputlisting[firstline=112,lastline=117,firstnumber=112]{../src/experiment.py}
\begin{longtable}{p{.14\textwidth}p{.76\textwidth}}
\toprule Lines & Exact operation\\\midrule\endfirsthead
\toprule Lines & Exact operation\\\midrule\endhead
112--113 & Helper function and docstring: select field occupations with one photon.\\
114--115 & List comprehension iterates tuple index $i$; the tuple itself is unused here.\\
116 & Require $\sum_{m<M}n_m(i)=1$. Thus there is exactly one occupied waveguide mode.\\
117 & Return matching tuple indices as an integer NumPy array, suitable for selecting columns of $C_g$.\\
\bottomrule
\end{longtable}

\subsection{Probabilities, direction and photon-count sectors}
\lstinputlisting[firstline=119,lastline=140,firstnumber=119]{../src/experiment.py}
Let $C_{j,i,s}$ be the returned amplitudes. The exact dictionary formulas are
\begin{align}
\mathcal N_j&=\sum_i\bigl(|C_{j,i,g}|^2+|C_{j,i,e}|^2\bigr),\\
P_{e,j}&=\sum_i|C_{j,i,e}|^2,\\
T_{1,j}&=\sum_{m:k_m>0}|C_{j,i(1_m),g}|^2,\qquad
R_{1,j}=\sum_{m:k_m<0}|C_{j,i(1_m),g}|^2,\\
P_{n,j}&=\sum_{i:\sum_{m<M}n_m(i)=n}
  \bigl(|C_{j,i,g}|^2+|C_{j,i,e}|^2\bigr).
\label{eq:obs}
\end{align}
\begin{longtable}{p{.14\textwidth}p{.76\textwidth}}
\toprule Lines & Exact operation / equation\\\midrule\endfirsthead
\toprule Lines & Exact operation / equation\\\midrule\endhead
119--121 & Reject measurement before \code{propagate_state}; coefficient arrays would still be \code{None}.\\
122--123 & \code{np.atleast_2d} turns final-only shape $(B,)$ into $(1,B)$; stored trajectory $(J,B)$ is unchanged. Both branches can now use \code{axis=1}.\\
124 & Elementwise $P_{j,i}=|C_{j,i,g}|^2+|C_{j,i,e}|^2$. \code{np.abs} is modulus, not real part.\\
125 & Get tuple indices with one waveguide photon.\\
126--127 & For each such tuple, \code{np.argmax(field occupations)} returns its unique occupied mode $m$ because total field count is exactly one.\\
128 & Keep one-photon tuple indices whose occupied $k_m>0$; right-moving set.\\
129 & Keep corresponding $k_m<0$ indices; left-moving set. A possible $k_m=0$ belongs to neither.\\
131--132 & Create output dictionary; time vector is all $t_j$ if stored or only final $t_{J-1}$ otherwise.\\
133 & Sum $P_{j,i}$ over \code{axis=1} (basis axis) to get $\mathcal N_j$.\\
134 & Sum excited amplitudes over basis axis to get $P_{e,j}$.\\
135 & Select ground-TLS amplitudes in right-moving one-photon tuples, square modulus and sum: $T_{1,j}$.\\
136 & Same for left-moving tuples: $R_{1,j}$.\\
137 & Close dictionary literal.\\
138 & Loop $n=0,\ldots,\max_i\sum_{m<M}n_m(i)$. For \code{full}, the largest $n$ can be $MN$, not $N$.\\
139 & Boolean mask \code{photon_numbers==n} selects occupation columns with $n$ waveguide photons. Sum their $P_{j,i}$ to get $P_{n,j}$. TLS excitation is not included in $n$.\\
140 & Return and retain dictionary of time arrays.\\
\bottomrule
\end{longtable}
$T_1+R_1$ need not equal one: TLS-excited and multi-photon sectors can retain
probability. These are finite-time sector probabilities; interpreting them as
asymptotic scattering requires the outgoing packet to have cleared the TLS
without periodic wraparound.

\subsection{Position-space one-photon amplitude}
\lstinputlisting[firstline=142,lastline=152,firstnumber=142]{../src/experiment.py}
\begin{equation}
\phi_{1,g}(x,t_j)=\frac1{\sqrt L}
\sum_{m=0}^{M-1}
C_{j,i(1_m),g}\,e^{\ii k_mx}.
\label{eq:phi}
\end{equation}
\begin{longtable}{p{.14\textwidth}p{.76\textwidth}}
\toprule Lines & Exact operation\\\midrule\endfirsthead
\toprule Lines & Exact operation\\\midrule\endhead
142--143 & Method signature and picture convention in docstring.\\
144--145 & Reject call before amplitudes exist.\\
146 & With full storage choose row $j=\texttt{i}$; with final-only storage use its sole vector, so argument \code{i} has no effect in that branch.\\
147 & Select one-waveguide-photon tuple indices.\\
148--149 & For each tuple find its unique occupied waveguide mode $m$.\\
150 & Convert requested positions $x_\ell$ to NumPy array.\\
151--152 & \code{np.outer(x_tab,k_tab[mode_index])} makes matrix $X_{\ell m}=x_\ell k_m$; exponentiation gives $e^{\ii X_{\ell m}}$; \code{@} multiplies by coefficients $C_{j,i(1_m),g}$; divide by $\sqrt L$. Output is $\phi_{1,g}(x_\ell,t_j)$ for each position.\\
\bottomrule
\end{longtable}
The solver already evolves Schr\"odinger-picture coefficients, so Eq.~\eqref{eq:phi}
needs no additional factor $e^{-\ii|k_m|t}$.

\section{Comparing two bases: \code{src/fidelities.py}}
For two normalized global kets $A,B$, three squared fidelities are returned:
\begin{equation}
F_{\rm state}=|\langle\psi_A|\psi_B\rangle|^2,\quad
F_{\rm atom}=\left[\Tr\sqrt{\sqrt{\rho_a^A}\rho_a^B
  \sqrt{\rho_a^A}}\right]^2,\quad
F_{\rm photon}=\left[\Tr\sqrt{\sqrt{\rho_f^A}\rho_f^B
  \sqrt{\rho_f^A}}\right]^2.
\label{eq:fids}
\end{equation}
The reduced matrices $\rho_a$ and $\rho_f$ mean partial traces over field and
TLS, respectively. They are \emph{diagnostics} derived from ket coefficients;
the global evolution remains closed and pure. The code computes $F_{\rm photon}$
without allocating a large field density matrix.

\subsection{Occupation-key map and TLS reduction}
\lstinputlisting[firstline=7,lastline=16,firstnumber=7]{../src/fidelities.py}
\begin{longtable}{p{.14\textwidth}p{.76\textwidth}}
\toprule Lines & Exact operation / equation\\\midrule\endfirsthead
\toprule Lines & Exact operation / equation\\\midrule\endhead
7--8 & Define the helper and make an empty dictionary: occupation key $\mapsto$ the two amplitudes $(C_g,C_e)$.\\
9 & Visit each occupation $\mathbf n_i$ in this experiment's own basis ordering.\\
10 & Save two-component vector \code{vector[2*i:2*i+2]} under the field occupation tuple, using Eq.~\eqref{eq:slice}. This changes labels without changing amplitudes.\\
11 & Return occupation-to-amplitude dictionary.\\
14--15 & Define the TLS reduction helper and reshape a length-$d$ ket into $B\times2$ coefficient matrix $C_{is}$; \code{-1} asks NumPy to infer $B=d/2$.\\
16 & \code{coefficients.T @ coefficients.conj()} yields $(\rho_a)_{ss'}=\sum_i C_{is}C^*_{is'}$. It is $2\times2$, Hermitian, trace one for a normalized ket. \code{.T} transposes; \code{.conj()} conjugates entrywise.\\
\bottomrule
\end{longtable}
The ordering of raw vectors differs among truncations. Comparing entry number
$q$ across two vectors would generally compare different physical occupations.
The dictionary at line 10 makes equality of keys, rather than equality of
integer positions, the criterion for a physical overlap.

\subsection{Preconditions, normalization and the $2\times2$ matrix $G$}
\lstinputlisting[firstline=19,lastline=44,firstnumber=19]{../src/fidelities.py}
In a common field space, decompose
\begin{equation}
\ket{\psi_A}=\ket{\phi_g^A}\ket g+\ket{\phi_e^A}\ket e,
\qquad
\ket{\psi_B}=\ket{\phi_g^B}\ket g+\ket{\phi_e^B}\ket e.
\end{equation}
Define the $2\times2$ cross-overlap matrix
\begin{equation}
G_{ss'}=\langle\phi_s^A|\phi_{s'}^B\rangle
=\sum_{\mathbf n\ \text{common}}
 C_{A,\mathbf n,s}^*C_{B,\mathbf n,s'},
\quad s,s'\in\{g,e\}.
\label{eq:G}
\end{equation}
An occupation absent from one basis has zero amplitude after embedding, hence
only keys in the intersection contribute.
\begin{longtable}{p{.14\textwidth}p{.76\textwidth}}
\toprule Lines & Exact operation / equation\\\midrule\endfirsthead
\toprule Lines & Exact operation / equation\\\midrule\endhead
19--20 & Function compares two completed experiments; docstring specifies squared fidelities.\\
21--22 & Require stored trajectories, since the function returns values at every $t_j$. Final-only objects do not have those lists.\\
23--24 & Require exactly equal signed arrays $k_m$ in the same order. A tuple entry $n_m$ must denote the same physical mode in both experiments.\\
25--26 & Require exactly equal output-time arrays; \code{zip} below then pairs states at the same time.\\
28 & Initialize three Python lists to collect scalar fidelity values over time.\\
29 & \code{zip} pairs $A$ and $B$ QuTiP kets at corresponding output indices.\\
30 & \code{ket.full()} is a two-dimensional column array; \code{[:,0]} yields a one-dimensional vector.\\
31--32 & Divide each vector by its Euclidean norm. Fidelity is defined for normalized states; this removes small ODE norm drift \emph{for metrics only}. Observable \code{norm} remains raw.\\
33--34 & Convert each vector into occupation-key dictionary with aligned TLS pairs.\\
35 & Allocate complex $2\times2$ matrix $G=0$.\\
36 & Python \code{a.keys() & b.keys()} is set intersection: iterate occupations physically present in both bases.\\
37 & \code{np.outer(a[key].conj(),b[key])} has entry $C^*_{A,\mathbf n,s}C_{B,\mathbf n,s'}$. Sum over keys yields Eq.~\eqref{eq:G}.\\
38 & $\Tr G=\sum_s\langle\phi_s^A|\phi_s^B\rangle=\langle\psi_A|\psi_B\rangle$, so append $|\Tr G|^2=F_{\rm state}$.\\
39 & \code{np.linalg.svd(G,compute_uv=False)} gives the two singular values $\sigma_1,\sigma_2$. Append $(\sigma_1+\sigma_2)^2=F_{\rm photon}$, the exact Uhlmann fidelity between the reduced field states.\\
40--41 & Construct each $2\times2$ TLS reduced matrix using line 16, wrapped for QuTiP.\\
42 & QuTiP returns root fidelity $\Tr\sqrt{\sqrt{\rho_a^A}\rho_a^B\sqrt{\rho_a^A}}$; square it to match Eq.~\eqref{eq:fids}.\\
43--44 & Convert lists to NumPy arrays of length $J$ and return dictionary with keys \code{state}, \code{atom}, \code{photon}.\\
\bottomrule
\end{longtable}
Why line 39 is exact: $\rho_f^A=
\ket{\phi_g^A}\bra{\phi_g^A}+\ket{\phi_e^A}\bra{\phi_e^A}$
has rank at most two, and likewise for $B$. The nonzero singular values needed
for Uhlmann fidelity are those of the $2\times2$ cross Gram matrix $G$.
Thus no $B\times B$ density matrix or $d^2$ global density matrix is needed.
The ratio $F_{\rm state}/F_{\rm atom}$ used elsewhere is not this fidelity.

\section{Script entry point: \code{experiments/scattering.py}}
\lstinputlisting[firstline=1,lastline=46,firstnumber=1]{../experiments/scattering.py}
\begin{longtable}{p{.16\textwidth}p{.74\textwidth}}
\toprule Lines & Exact operation\\\midrule\endfirsthead
\toprule Lines & Exact operation\\\midrule\endhead
1--7 & Docstring and imports: \code{csv} writes optional output; \code{sys} and \code{Path} locate modules/files; \code{np} is only used by the currently commented example.\\
9 & \code{Path(__file__).resolve().parents[1]} is the \code{codex/code} project root.\\
10 & Insert project root at the front of Python's import path.\\
12--13 & Import configuration and experiment classes after path insertion.\\
16--22 & Four required dictionaries; optional $N$, truncation, RWA, state storage, CSV writing and progress. Defaults: $N=3$, \code{full+totalcap}, \code{RWA=False}, \code{store_state=True}, no file, no progress bar.\\
24 & Docstring: return an \code{Experiment} object.\\
26--33 & Construct \code{ExperimentConfig} with explicit keyword assignments. \code{store_results} and \code{progress} control I/O, not physics, and are not in the config.\\
35 & Constructor builds grid, basis, $H_0,V_1,H$, and initial ket.\\
36 & QuTiP propagates Eq.~\eqref{eq:ode}; save amplitudes.\\
37 & Compute Eq.~\eqref{eq:obs}; dictionary also stored on the experiment.\\
39 & Enter optional CSV branch only if \code{store_results=True}.\\
40 & Set output path \code{results/csv_files/scattering.csv} under project root.\\
41 & Create parent directories as needed; existing directories are accepted.\\
42 & Open file for writing; existing CSV at this path is replaced. \code{newline=""} delegates line endings to CSV module.\\
43 & Create a CSV writer.\\
44 & Write metric names as header in dictionary insertion order.\\
45 & \code{zip(*observables.values())} transposes metric arrays: row $j$ contains all values at $t_j$.\\
46 & Return complete experiment for notebook inspection.\\
\bottomrule
\end{longtable}
Lines 49--62 are presently \emph{comments}. Running
\code{python3 experiments/scattering.py} currently defines the function and
exits without a simulation. The README's run command describes an earlier
version; the active calls are in the notebook.

\section{Notebook \code{notebooks/scattering.ipynb}: current cells}
Cell numbers below start at zero, matching their JSON order. Markdown cells
0, 5 and 7 are text only. Stored notebook \emph{outputs} are historical; a
source cell may have changed since they were generated. Outputs of comparison
cells 6 and 8 were cleared when the basis options changed.

\subsection{Cell 1: imports and module lookup}
\begin{lstlisting}
import sys
from pathlib import Path
import numpy as np
import matplotlib.pyplot as plt
project_root = Path.cwd().parent
sys.path.insert(0, str(project_root))
from src.experiment import resource_estimation
from src.fidelities import fidelities_over_time
from experiments.scattering import run_scattering
\end{lstlisting}
\begin{longtable}{p{.25\textwidth}p{.65\textwidth}}
\toprule Statement & Exact role\\\midrule\endfirsthead
\toprule Statement & Exact role\\\midrule\endhead
\code{import sys} & Access module search path.\\
\code{Path} & \code{from pathlib import Path}: filesystem path object.\\
\code{import numpy as np} & Short name for arrays, $\pi$, means, nearest-index selection and geometric grids.\\
\code{import matplotlib.pyplot as plt} & Short name for plotting.\\
\code{Path.cwd().parent} & Assumes the working directory is \code{codex/code/notebooks}; parent is the project root.\\
\code{sys.path.insert} & Put that root first in Python's import search list.\\
\code{resource_estimation} & Import $M,d,J$ and memory estimator.\\
\code{fidelities_over_time} & Import exact comparison of two trajectories.\\
\code{run_scattering} & Import construction, propagation and measurement wrapper.\\
\bottomrule
\end{longtable}

\subsection{Cell 2: physical inputs and advisory count}
\begin{lstlisting}
param_photon = {'k_0': 1.0, 'sigma_k': 0.10, 'x_0': -12.0}
param_atom = {'omega_0': 1.0, 'D': 0.45, 'x_tls': 0.0,
              'L': 20*np.pi, 'coupling': 'sqrt'}
param_time_evol = {'T': 28.0, 'dt': 0.1}
cutoffs = {'ir_cutoff': 0.8, 'uv_cutoff': 1.2}
n_max = 3
truncation = 'full'
estimate = resource_estimation(param_atom, param_time_evol, cutoffs,
                               n_max=n_max, truncation=truncation)
\end{lstlisting}
\begin{longtable}{p{.25\textwidth}p{.65\textwidth}}
\toprule Statement & Exact meaning / consequence\\\midrule\endfirsthead
\toprule Statement & Exact meaning / consequence\\\midrule\endhead
\code{param_photon} & Eq.~\eqref{eq:packet}: $k_0=1$, $\sigma_k=0.1$, $x_0=-12$. Absent \code{state} means one photon; absent \code{right_moving_only} means $k>0$ only.\\
\code{param_atom} & $\omega_0=1,D=0.45,x_a=0,L=20\pi$, $f_m=\sqrt{|k_m|}$. Absent \code{initial_state} means ground TLS.\\
\code{param_time_evol} & $T=28,\Delta t=0.1$; $J=281$ saved times $0,0.1,\ldots,28$.\\
\code{cutoffs} & Eq.~\eqref{eq:grid}: $M=10$ signed modes, five per direction.\\
\code{n_max} & $N=3$.\\
\code{truncation} & Current value \code{'full'} independently caps ten modes at 3: $d=2(4)^{10}=2{,}097{,}152$. Function default \code{full+totalcap} would give $2\binom{13}{3}=572$.\\
\code{estimate} & Prints approximately 8.78 GiB for stored complex amplitudes and \code{feasible=False}. It does not prevent execution.\\
\bottomrule
\end{longtable}

\subsection{Cell 3: the simulation actually requested}
\begin{lstlisting}
experiment = run_scattering(param_photon, param_atom,
                            param_time_evol, cutoffs,
                            n_max=n_max, truncation=truncation,
                            RWA=True, progress=True)
for name, values in experiment.observables.items():
    print(name, values[-1])
\end{lstlisting}
\begin{longtable}{p{.25\textwidth}p{.65\textwidth}}
\toprule Statement & Exact meaning / consequence\\\midrule\endfirsthead
\toprule Statement & Exact meaning / consequence\\\midrule\endhead
\code{run_scattering(...)} & Construct and solve with cell-2 inputs. This cell does \emph{not} test \code{estimate['feasible']}; it can begin a very large allocation.\\
\code{RWA=True} & Override default \code{False}; this cell uses only energy-exchange transitions. It is a rotating-wave calculation despite the notebook title.\\
\code{progress=True} & Enable progress bar, with no change in physics.\\
\code{for name,values} & Iterate output dictionary, one metric array at a time.\\
\code{values[-1]} & Python index $-1$ means the last \emph{time} sample, here $t=28$.\\
\bottomrule
\end{longtable}

\subsection{Cell 4: spatial and time plots}
\begin{lstlisting}
if experiment is not None:
    time = experiment.observables['time']
    x_tab = np.linspace(-param_atom['L']/2, param_atom['L']/2,
                        500, endpoint=False)
    t_hit = param_atom['x_tls'] - param_photon['x_0']
    snapshots = [0, np.argmin(np.abs(time - t_hit)), len(time)-1]
    fig, axes = plt.subplots(1, 2, figsize=(11, 4))
    for j in snapshots:
        phi_x = experiment.one_photon_wavefunction(j, x_tab)
        axes[0].plot(x_tab, np.abs(phi_x)**2, label=f't = {time[j]:.1f}')
    axes[0].axvline(param_atom['x_tls'], color='black', ls=':', label='TLS')
    axes[0].set(xlabel='x', ylabel=r'$|\phi_{1,g}(x,t)|^2$')
    axes[0].legend()
    for name in ('p_transmitted','p_reflected','p_excited','norm'):
        axes[1].plot(time, experiment.observables[name], label=name)
    axes[1].set(xlabel='t', ylabel='population', ylim=(0, 1.05))
    axes[1].legend()
    fig.tight_layout()
    plt.show()
\end{lstlisting}
\begin{longtable}{p{.28\textwidth}p{.62\textwidth}}
\toprule Statement & Mathematical meaning\\\midrule\endfirsthead
\toprule Statement & Mathematical meaning\\\midrule\endhead
\code{if experiment is not None} & Plot if a previous cell assigned an experiment; a failed previous cell can leave a stale notebook variable.\\
\code{time} & Read solver output array $(t_j)$.\\
\code{x_tab} & Sample 500 points on $[-L/2,L/2)$; \code{endpoint=False} avoids duplicating periodic endpoints.\\
\code{t_hit} & $x_a-x_0=12$, rough flight time at speed $c=1$, not a computed interaction time.\\
\code{snapshots} & Indices $0$, nearest saved time to $12$ via \code{argmin(abs(...))}, and $J-1$.\\
\code{plt.subplots} & Allocate one figure and two axes.\\
\code{for j}, \code{phi_x} & Compute Eq.~\eqref{eq:phi} at three chosen $t_j$.\\
\code{axes[0].plot} & Plot $|\phi_{1,g}(x,t_j)|^2$ and formatted time label.\\
\code{axvline}, first \code{set}, \code{legend} & Mark TLS position $x_a$; label spatial axes/curves. No new physics is calculated.\\
\code{for name}, \code{axes[1].plot} & Plot $T_1(t_j),R_1(t_j),P_e(t_j),\mathcal N(t_j)$ from Eq.~\eqref{eq:obs}.\\
second \code{set}, \code{legend} & Label time plot and restrict visible vertical range to $(0,1.05)$.\\
\code{tight_layout}, \code{show} & Arrange and display; no raw data file saved.\\
\bottomrule
\end{longtable}

\subsection{Cell 6: two truncations and three fidelities}
\begin{lstlisting}
schemes = ('full+totalcap', 'truncated')
runs = {}
for scheme in schemes:
    estimate = resource_estimation(param_atom, param_time_evol,
                                   cutoffs, n_max=n_max, truncation=scheme)
    if estimate['feasible']:
        runs[scheme] = run_scattering(param_photon, param_atom,
                                     param_time_evol, cutoffs,
                                     n_max=n_max, truncation=scheme)
reference = runs.get('full+totalcap')
if reference is not None:
    fig, axes = plt.subplots(1, 3, figsize=(13, 3.5))
    for scheme in ('truncated',):
        if scheme in runs:
            F = fidelities_over_time(runs[scheme], reference)
            for ax, name in zip(axes, ('state', 'atom', 'photon')):
                ax.plot(reference.times, F[name], label=scheme)
                ax.set(xlabel='t', ylabel=f'{name} fidelity',
                       ylim=(0, 1.05))
                ax.set_title(name)
            print(scheme, 'initial:', F['state'][0],
                  'time mean:', {key: np.mean(value)
                                for key, value in F.items()})
    axes[0].legend()
    fig.tight_layout()
    plt.show()
\end{lstlisting}
\begin{longtable}{p{.28\textwidth}p{.62\textwidth}}
\toprule Statement & Exact operation\\\midrule\endfirsthead
\toprule Statement & Exact operation\\\midrule\endhead
\code{schemes}, \code{runs} & List two finite bases; initialize name-to-experiment dictionary.\\
\code{for scheme}, \code{estimate} & Recalculate $d,J$ for each basis with the same physical dictionaries.\\
\code{if estimate['feasible']} & This cell \emph{does} obey the advisory threshold.\\
\code{runs[scheme]} & Solve feasible cases. No \code{RWA} keyword is supplied, so wrapper default \code{False} applies: these runs include counter-rotating terms, unlike cell 3.\\
\code{reference=...} & Retrieve total-cap result, or \code{None} if skipped.\\
\code{if reference...}, \code{subplots} & Make three axes only if reference exists.\\
\code{for scheme}, \code{if scheme in runs} & Consider each available candidate against that same finite reference.\\
\code{F=...} & Compute all $F_\mu(t_j)$ in Eq.~\eqref{eq:fids}; requires common $k$ array and $t$ array.\\
\code{zip(axes,...)} & Pair axis 0/1/2 with names \code{state}/\code{atom}/\code{photon}.\\
\code{ax.plot}, \code{ax.set}, \code{ax.set_title} & Plot $F_\mu(t_j)$, label axes and use visible range $(0,1.05)$.\\
\code{F['state'][0]} & Initial global fidelity $F_{\rm state}(0)$.\\
\code{np.mean(value)} & Arithmetic sampled-time mean $J^{-1}\sum_{j=0}^{J-1}F_\mu(t_j)$ for each metric.\\
\code{legend}, \code{tight_layout}, \code{show} & Display comparison; no file saved.\\
\bottomrule
\end{longtable}

\subsection{Cell 8: coupling sweep, currently enabled}
\begin{lstlisting}
RUN_SWEEP = True
if RUN_SWEEP:
    D_values = np.geomspace(0.02, 0.45, 6)
    sweep = {scheme: [] for scheme in ('truncated',)}
    for D in D_values:
        atom_here = dict(param_atom, D=float(D))
        ref = run_scattering(param_photon, atom_here, param_time_evol,
                             cutoffs, n_max=n_max, truncation='full+totalcap')
        for scheme in sweep:
            candidate = run_scattering(param_photon, atom_here,
                                       param_time_evol, cutoffs,
                                       n_max=n_max, truncation=scheme)
            F = fidelities_over_time(candidate, ref)
            sweep[scheme].append(np.mean(F['state']))
    for scheme, values in sweep.items():
        plt.semilogx(D_values, values, 'o-', label=scheme)
    plt.xlabel('D')
    plt.ylabel('time-mean full-state fidelity')
    plt.legend()
    plt.show()
\end{lstlisting}
\begin{longtable}{p{.28\textwidth}p{.62\textwidth}}
\toprule Statement & Exact operation\\\midrule\endfirsthead
\toprule Statement & Exact operation\\\midrule\endhead
\code{RUN_SWEEP=True}, \code{if} & Current value enters the sweep branch. Notebook Markdown says to enable it, but it is already enabled.\\
\code{D_values} & Six geometrically spaced $D$: equal spacing in $\log D$ between $0.02$ and $0.45$.\\
\code{sweep} & One empty output list keyed by candidate basis.\\
\code{for D}, \code{atom_here} & For each $D$, shallow-copy atom dictionary and replace only $D$; original \code{param_atom} remains unchanged.\\
\code{ref} & One total-cap simulation at this $D$ and default \code{RWA=False}.\\
\code{for scheme}, \code{candidate} & One candidate simulation per $D$: $6(1+1)=12$ solver calls total, each rebuilding a matrix.\\
\code{F}, \code{append} & Compute full-state fidelity trajectory and append its time mean.\\
\code{for scheme,values}, \code{semilogx} & Plot six means against logarithmic $D$ axis.\\
\code{xlabel}, \code{ylabel}, \code{legend}, \code{show} & Label/display; no raw sweep array saved.\\
\bottomrule
\end{longtable}

\section{A complete small-vector example}
Take $M=2,N=1$, \code{truncated}, normalized packet
$c=(c_0,c_1)$, and TLS spinor $\eta=(\eta_g,\eta_e)$.
The basis constructor stores occupation tuples in order
\begin{align*}
\texttt{basis.states}&=[(0,0),(1,0),(0,1)],\\
\texttt{basis.index}[(0,0)]&=0,\quad
\texttt{basis.index}[(1,0)]=1,\quad
\texttt{basis.index}[(0,1)]=2.
\end{align*}
Consequently the full ket vector has six cells:
\[
\begin{array}{c|cccccc}
q&0&1&2&3&4&5\\
\hline
\text{ket}&\ket{0,0;g}&\ket{0,0;e}&\ket{1,0;g}&
\ket{1,0;e}&\ket{0,1;g}&\ket{0,1;e}\\
v_q&0&0&c_0\eta_g&c_0\eta_e&c_1\eta_g&c_1\eta_e
\end{array}
\]
The \code{for m, amplitude in enumerate(packet)} loop writes
\[
\begin{aligned}
m=0:\quad &i=1,\quad
  v[2:4]\mathrel{+}=c_0(\eta_g,\eta_e),\\
m=1:\quad &i=2,\quad
  v[4:6]\mathrel{+}=c_1(\eta_g,\eta_e).
\end{aligned}
\]
The vacuum cells remain zero because a one-photon number state has no vacuum
component. Its norm is
\[
\|v\|^2
=|c_0|^2(|\eta_g|^2+|\eta_e|^2)
 +|c_1|^2(|\eta_g|^2+|\eta_e|^2)
=1.
\]
This is the literal contents of the array after line 79 of \code{src/states.py};
the final normalization at line 93 leaves it unchanged up to floating-point
roundoff.

\section{Imports, package markers, and files without numerical steps}
The executable listings above start at the classes/functions. The leading
imports select the following exact implementations:
\begin{center}\small
\begin{longtable}{p{.29\textwidth}p{.61\textwidth}}
\toprule Source lines & Purpose\\\midrule\endfirsthead
\toprule Source lines & Purpose\\\midrule\endhead
\code{src/states.py:1--7} & Module docstring; \code{itertools.combinations_with_replacement} and \code{product} enumerate occupations; \code{math.factorial} supplies $n!$ in Eq.~\eqref{eq:coherent}; NumPy makes arrays; QuTiP wraps final ket.\\
\code{src/hamiltonians.py:1--5} & Module docstring; NumPy makes $f_m$ and phases; QuTiP wraps $H$; SciPy \code{coo_matrix} stores triplets and \code{diags} stores $H_0$.\\
\code{src/experiment.py:1--9} & Module docstring; \code{comb} gives binomial count, \code{pi} the momentum spacing, NumPy arrays, QuTiP solver, and imports the local basis/Hamiltonian functions.\\
\code{src/fidelities.py:1--4} & Module docstring; NumPy computes overlaps/SVD; QuTiP computes TLS Uhlmann fidelity.\\
\code{src/__init__.py}, \code{experiments/__init__.py} & Empty package-marker files. No numerical statements.\\
\code{requirements.txt} & Declares NumPy, SciPy, QuTiP, Matplotlib dependency ranges; it supplies no physical parameter.\\
\code{README.md} & Short usage overview. Its script command is currently stale because the script's example is commented out.\\
\code{results/} & Optional CSV output location; no result is written unless \code{store_results=True}.\\
\bottomrule
\end{longtable}
\end{center}

\section{What the equations certify, and what they do not}
\begin{center}\small
\begin{longtable}{p{.40\textwidth}p{.48\textwidth}}
\toprule Exact finite-code identity & Scope\\\midrule\endfirsthead
\toprule Exact finite-code identity & Scope\\\midrule\endhead
$H=H_0+D V_1$, Eqs.~\eqref{eq:HD}--\eqref{eq:V} & Sparse projected Hamiltonian for the chosen $k$ and Fock basis.\\
$H^\dagger=H$ for real $D,L>0,\omega_0$ & Down/up entries are conjugate; exact finite evolution is unitary. A high norm cannot establish basis convergence.\\
Default $v_0=\sum_m c_m\ket{1_m;g}$ & A right-moving, one-photon preparation. Coherent preparation is normalized \emph{after} projection.\\
Eq.~\eqref{eq:obs} & $T_1,R_1$ are ground-TLS, one-photon sector weights at finite time; they are not automatically asymptotic scattering coefficients.\\
$F_{\rm photon}=(\sigma_1(G)+\sigma_2(G))^2$ & Exact field reduced-state fidelity on matching $k,t$ grids, without a field density matrix.\\
\code{store_state=False} & Only final coefficients/observables are retained; time-dependent fidelity function requires full stored trajectories.\\
\bottomrule
\end{longtable}
\end{center}
Numerical convergence has independent axes: photon cap $N$, mode band,
periodic length $L$, and ODE tolerances. Changing one while fixing the others
tests that axis only. The current notebook cell 3 uses RWA and an enormous
\code{full} basis; cells 6 and 8 use no-RWA by default. Their curves do not
represent one single parameter regime unless those switches are aligned.

\end{document}
